\section{Conclusion}
Under the evaluated full-load sequential-disjoint topology, increasing $N$ increases the number of simultaneous directed links and drives the shared-resource case into a severe interference-limited regime. Resource partitioning and the experimental desired/interferer advantage move the mean-based operating region substantially. The robustness campaign also shows why this result must be topology-conditioned: at $N=100$, greedy short-link pairing raises shared-resource first-TX success probability from 0.0128 to 0.4766 by shortening desired links. Geometry-aware reuse and adaptive MCS provide additional gains, while the tested noise-bandwidth convention has negligible effect. The supported design message is therefore conditional but actionable: topology, reuse, spatial selectivity, and link adaptation should be evaluated jointly rather than inferring swarm capacity from $N$ alone.

\section*{Acknowledgment}
OpenAI ChatGPT was used as an editing assistant for language, compression, LaTeX formatting, and bibliographic organization. It was not used to generate simulation data. The author retains responsibility for the scientific content, source verification, and final manuscript.
